% ===============================================================================
% = LaTeX Beamer Template des Arbeitsbereichs Sicherheit in verteilten Systemem
% = (c) 2016 Prof. Dr. Hannes Federrath, Uni Hamburg, Fachbereich Informatik
% = https://svs.informatik.uni-hamburg.de
% =
% = Weitgehend in Übereinstimmung mit dem Corporate Design 2016 der UHH:
% = https://www.uni-hamburg.de/beschaeftigtenportal/services/oeffentlichkeitsarbeit/corporate-design.html
% = 
% ===============================================================================
%
\documentclass[t,aspectratio=169]{beamer} 
% Option t              Place text of slides at the (vertical) top of the slides.
% Option handout        Ein PDF ohne Pausen und Overlayeffekte erzeugen.
% Option aspectratio=43 169 => 16:9, 1610 => 16:10, 43 => 4:3
\usepackage[utf8]{inputenc}
\usepackage[ngerman]{babel}
\usepackage{graphicx,xcolor}
\usepackage[T1]{fontenc} % 8-Bit-Zeichen; ermöglicht korrektes Kopieren von Umlauten aus dem pdf 
\bibliography{bib}
\bibliographystyle{ieee}
\usepackage{biblatex}
\addbibresource{bib.bib}

\newcommand{\numUsers}{n}
\newcommand{\numElements}{k}

% SVS-Theme benutzen
\usetheme{svs2021}

% =============================
% = Ab hier Inhalte ändern... 
% =============================

\title{Security and Privacy in Federated Learning}
% \subtitle{Ein Vorschlag}
\author[]{Jonas Müller, Rebekka Schnoor}
% \institute[Uni Hamburg]{Universität Hamburg · Fachbereich Informatik}
\date{21.05.2026} % Ausblenden geht mit leerem Inhalt: \date{}
% \titlegraphic{\uhhlogo} % Logo für Titelseite, UHH-Logo ist Standard
\eventlocation{Master Project Seminar Data Protection and Data Security, Hamburg, 21$^{st}$ of May 2026} % rechts oben auf Titelseite

\begin{document}

\begin{frame}
	% Die Titelseite erscheint nach erneutem Übersetzen korrekt.
	\maketitle
\end{frame}

\begin{frame}
	\frametitle{Agenda}
	% Die Gliederung erscheint nach erneutem Übersetzen korrekt.
	% Die Gliederung wird im Dokument fortlaufend durch \section{} und \subsection{} definiert. Die Folien selbst haben keinen Einfluss auf die Gliederung.
	\tableofcontents
	% Hinweis: Das pandoc-Template benutzt \tableofcontents[hideallsubsections] zur Erzeugung einer Gliederung. Daher werden in der Markdown-Version nur die Überschriften der ersten Ebene (#) im Inhaltsverzeichnis angezeigt. Die normalerweise bei Verwendung von pandoc erzeugten Zwischenfolien für jede Überschrift der ersten Ebene werden im Theme unterdrückt. 
\end{frame}

\section{Secure Aggregation}
\subsection{Motivation}
\subsection{Threat Models}
\subsection{Protocol Examples}

\begin{frame}{Secure Aggregation}
\begin{alertblock}{Motivation}
    Even if the datasets stay with the clients, shared gradients can disclose information and enable inference attacks for both centralized and distributed FL.
\end{alertblock}
\pause

\textbf{Threat Models}
\begin{itemize}
    \item Honest-but-curious
    \begin{itemize}
        \item[] The participant follows the protocol honestly but tries to learn as much as possible from the received messages
    \end{itemize}
    \pause
    \item Malicious
    \begin{itemize}
        \item[] The participant can lie or break the protocol
    \end{itemize}
    \pause
    \item Malicious + (limited) access to memory
    \begin{itemize}
        \item[] The malicious participant additionally has access to some of the other participant's local memory (i.e. learns their secrets)
    \end{itemize}
\end{itemize} 
\end{frame}

\begin{frame}{Secure Aggregation}
    
\begin{alertblock}{Motivation}
    Even if the datasets stay with the clients, shared gradients can disclose information and enable inference attacks for both centralized and distributed FL.
\end{alertblock}
\begin{columns}
\begin{column}{0.5\textwidth}
\textbf{Threat Models}
\begin{itemize}
    \item Honest-but-curious
    \item[] \textcolor{black}{Follow protocol but maximize learning}
    \item Malicious
    \item[] \textcolor{black}{Can break the protocol}
    \item Malicious + (limited) access to memory
    \item[] \textcolor{black}{+ Partial access to secrets}
\end{itemize} 
\end{column}
\begin{column}{0.5\textwidth}
    \textbf{Challenges}
    \begin{itemize}
        \item User dropout
        \item Limited bandwidth and/or computation power
        \item Transmission of high dimensional information vectors (gradients)
        \item Unauthenticated network model
    \end{itemize}
\end{column}
\end{columns}
\end{frame}

\begin{frame}{Secure Aggregation Protocol by Bonawitz et al. \cite{bonawitz2016practicalsecureaggregationfederated}}
\begin{columns}[T]
    
    \begin{column}{0.5 \textwidth}
        \textbf{Protocol Summary}
        \begin{itemize}
            \item Uses masking vectors to encrypt individual gradients that cancel out in the aggregation
            \item Robust to user dropout by using secret sharing and a second masking vector
            \item Secure against malicious adversaries with limited memory access by using OneTime Pads and a sufficiently high threshold for secret reconstruction
        \end{itemize}
    \end{column}
    \begin{column}{0.5\textwidth}
        \textbf{Costs} \\ for $n$ users and vector dimension $k$

          \begin{tabular}{ll}
            \\
            \hline
            \multicolumn{2}{l}{\textbf{computation}} \\
            \hline 
            User & $O(\numUsers^2 + \numElements\numUsers)$ \\
            Server & $O(\numElements\numUsers^2)$\\
            \hline
            \multicolumn{2}{l}{\textbf{communication}} \\
            \hline
            User & $O(\numUsers + \numElements)$ \\
            Server & $O(\numUsers^2 + \numElements\numUsers)$ \\
            \hline
            \multicolumn{2}{l}{\textbf{storage}} \\
            \hline
            User & $O(\numUsers + \numElements)$ \\
            Server & $O(\numUsers^2 + \numElements)$ \\
            
            \end{tabular}
            \label{table:costs}
    \end{column}
\end{columns}
\end{frame}

\begin{frame}{SAFELearn: Secure Aggregation for private Federated Learning \cite{cryptoeprint:2021/386}}

            \textbf{2 Main Variants:}\\
            \begin{itemize}
                \item Fully Homomorphic Encryption (FHE)
                \item [] \textcolor{gray}{Uses additively homomorphic encryption scheme}
                \item [] \textcolor{gray}{Encryption key is split among clients using secret sharing}
                \begin{enumerate}
                    \item Clients send encrypted gradients to the server
                    \item Server aggregates encrypted gradients and sends result to clients
                    \item Clients can then jointly decrypt the result
                \end{enumerate}
                \item Secure Multi-Party Computation (MPC) / Secure Two Party Computation (2PC)
                \begin{enumerate}
                    \item Clients split their updates among N non-colluding servers using secret sharing
                    \item Servers jointly aggregate and send shares back to clients 
                    \item Clients reconstruct the result
                \end{enumerate}
            \end{itemize}
\end{frame}

\begin{frame}{References}
\printbibliography
    
\end{frame}

%%%%%%%%%%%%%%%%%%%%% bonus slides %%%%%%%%%%%%%%%%%%%%%
\begin{frame}{Secure Aggregation Protocol by Bonawitz et al. \cite{bonawitz2016practicalsecureaggregationfederated}}
    \begin{columns}
        \begin{column}{0.5\textwidth}
            %protocol steps
            \begin{enumerate}
                \item Users pairwise agree on input perturbation pair $p_{u,v}, p_{v,u}$
                \item Each user masks their gradient using the sum of all perturbations known and sends the result to the Server
                \item The Server aggregates as usual because the perturbations are structured to cancel each other
            \end{enumerate}
        \end{column}
        \pause
        \begin{column}{0.5\textwidth}
            %up- and downsides
            \begin{itemize}
                \item[\textbf{\textcolor{green}{+}}] information theoretically secure through OneTime Pads, even for malicious adversaries with limited memory access
                \pause
                \item[\textbf{\textcolor{red}{!}}] strong assumptions:
                \begin{itemize}
                    \item[] pairwise secure channels with sufficient bandwidth for secret exchange  
                    \item[] protocol completion by all users
                \end{itemize}
            \end{itemize}
        \end{column}
    \end{columns}
\end{frame}

\begin{frame}{Secure Aggregation Protocol by Bonawitz et al. \cite{bonawitz2016practicalsecureaggregationfederated}}
    \begin{columns}
        \begin{column}{0.6\textwidth}
            %protocol steps
            \begin{enumerate}
            \setcounter{enumi}{-1}
                \item Each user creates a key pair and broadcasts their public key to all other users

                \item \textcolor{gray}{Perturbation vector agreement}
                \begin{enumerate}
                    \item Secret sharing: Users distribute perturbation vector among all other users with reconstruction threshold $t$
                \end{enumerate}
                
                \item \textcolor{gray}{Masking + sending of gradients}
                \begin{enumerate}
                    \item Only use perturbations shared between surviving users. Reveal secrets of dropped out users if at least $t$ users survived
                \end{enumerate}
                
                \item \textcolor{gray}{Aggregation by server
                }
            \end{enumerate}
        \end{column}
        \pause
        \begin{column}{0.4\textwidth}
            %up- and downsides
            \begin{itemize}
                \item[\textbf{\textcolor{green}{+}}] Robustness to user dropout
                \pause
                \item[\textbf{\textcolor{red}{!}}] Secret exposure in case of malicious server or delayed update sending; increased communication cost
            \end{itemize}
        \end{column}
    \end{columns}
\end{frame}

\begin{frame}{Secure Aggregation Protocol by Bonawitz et al. \cite{bonawitz2016practicalsecureaggregationfederated}}
    \begin{columns}
        \begin{column}{0.5\textwidth}
            %protocol steps
            \begin{enumerate}
            \setcounter{enumi}{-1}
                \item \textcolor{gray}{Key pair generation} + generation of masking vector $b$
                \item \textcolor{gray}{Perturbation vector agreement}
                \begin{enumerate}
                    \item \textcolor{gray}{Secret sharing} including $b$
                \end{enumerate}
                \item \textcolor{gray}{Masking + sending of gradients:} Use $b$ as second mask
                \begin{enumerate}
                    \item \textcolor{gray}{Reveal secrets of dropped out users if at least $t$ users survived:} Each user will only reveal a share of $p_u$ OR $b_u$
                \end{enumerate}
                \item \textcolor{gray}{Aggregation by server} 
                
            \end{enumerate}
        \end{column}
        \pause
        \begin{column}{0.5\textwidth}
            %up- and downsides
            \begin{itemize}
                \item[\textbf{\textcolor{green}{+}}] For $t > \frac{n}{2}$, all gradients stay masked by either $p$ or $b$; extend to stronger adversaries by increasing threshold
                \pause
                \item[\textbf{\textcolor{red}{!}}] Communication cost of $O(kn^2)$
            \end{itemize}
            \pause
            \begin{alertblock}{Reducing Complexity}
                    Share seeds for PRG instead of full vectors using DH Key Exchange via server
            \end{alertblock}
            \pause
            \begin{alertblock}{Getting rid of secured channel assumption}
                Setup secured channels via server-mediated DH  $\rightarrow$ adds PFS
            \end{alertblock}
        \end{column}
    \end{columns}
\end{frame}
\end{document}